% concatenated from Law_of_Learning_Dynamics.tex, preamble.tex, sec/4_Experiment.tex, sec/5_Conclusion.tex
%%%%%%%% ICML 2026 EXAMPLE LATEX SUBMISSION FILE %%%%%%%%%%%%%%%%%

\documentclass{article}

% Recommended, but optional, packages for figures and better typesetting:
\usepackage{microtype}
\usepackage{graphicx}
\usepackage{subcaption}
\usepackage{booktabs} % for professional tables
\usepackage{makecell} 
% hyperref makes hyperlinks in the resulting PDF.
% If your build breaks (sometimes temporarily if a hyperlink spans a page)
% please comment out the following usepackage line and replace
% \usepackage{icml2026} with \usepackage[nohyperref]{icml2026} above.
\usepackage{hyperref}
\usepackage{subcaption}

% Attempt to make hyperref and algorithmic work together better:
\newcommand{\theHalgorithm}{\arabic{algorithm}}

% Use the following line for the initial blind version submitted for review:
%\usepackage{icml2026}

% For preprint, use
\usepackage[preprint]{icml2026}

% If accepted, instead use the following line for the camera-ready submission:
% \usepackage[accepted]{icml2026}

\usepackage{amsmath}
\usepackage{amssymb}
\usepackage{mathtools}
\usepackage{amsthm}
\usepackage{bbm}


% if you use cleveref..
\usepackage[capitalize,noabbrev]{cleveref}

%%%%%%%%%%%%%%%%%%%%%%%%%%%%%%%%
% THEOREMS
%%%%%%%%%%%%%%%%%%%%%%%%%%%%%%%%
\theoremstyle{plain}
\newtheorem{theorem}{Theorem}[section]
\newtheorem{proposition}[theorem]{Proposition}
\newtheorem{lemma}[theorem]{Lemma}
\newtheorem{corollary}[theorem]{Corollary}
\theoremstyle{definition}
\newtheorem{definition}[theorem]{Definition}
\newtheorem{assumption}[theorem]{Assumption}
\theoremstyle{remark}
\newtheorem{remark}[theorem]{Remark}

%%%%%%%%%%%%%%%%%%%%%%%%%%%%%%%%
% NEWCOMMANDS in preamble.tex
%%%%%%%%%%%%%%%%%%%%%%%%%%%%%%%%
%% This file contains a number of tweaks that are typically applied to the main document.
%% They are not enabled by default, but can be enabled by uncommenting the relevant lines.


\newtheorem{thm}{Theorem}[section]
\newtheorem{prop}[thm]{Proposition}
\newtheorem{lem}[thm]{Lemma}
\newtheorem{rem}[thm]{Remark}
\newtheorem{cor}[thm]{Corollary}
\newtheorem{exr}[thm]{Exercise}
\newtheorem{PG}[thm]{Conjecture}
\newtheorem{defn}[thm]{Definition}

\newcommand{\red}[1]{{\color{red}#1}}
%\newcommand{\todo}[1]{{\color{red}#1}}
\newcommand{\TODO}[1]{\textbf{\color{red}[TODO: #1]}}

\newcommand{\N}{\mathbb{N}}
\newcommand{\Z}{\mathbb{Z}}
\newcommand{\R}{\mathbb{R}}
\newcommand{\U}{\mathbb{U}}
\newcommand{\Q}{\mathbb{Q}}
\newcommand{\C}{\mathbb{C}}
\newcommand{\bS}{\mathbb{S}}
\newcommand{\bi}{\mathbf{i}}
\newcommand{\bb}{\boldsymbol{b}}
\newcommand{\bv}{\mathbf{v}}
\newcommand{\bw}{\mathbf{w}}
\newcommand{\bx}{\boldsymbol{x}}
\newcommand{\WH}[1]{\widehat{#1}}
\newcommand{\CE}{\mathcal{E}}
\newcommand{\CF}{\mathrm{CF}}
\newcommand{\HF}{\mathrm{HF}}
\newcommand{\CO}{\mathcal{O}}
\newcommand{\CC}{\mathcal{C}}
\newcommand{\CL}{\mathcal{L}}
\newcommand{\cY}{\mathcal{Y}}
\newcommand{\bL}{\mathbb{L}}
\newcommand{\BL}{\mathbf{L}}
\newcommand{\A}{\mathbb{A}}
\newcommand{\CP}{\mathbb{C}P}
\newcommand{\bD}{\mathcal{D}}
\newcommand{\uF}{\underline{F}}
\newcommand{\cF}{\mathscr{F}}
\newcommand{\cP}{\mathcal{P}}
\newcommand{\cS}{\mathcal{S}}
\newcommand{\cT}{\mathcal{T}}
\newcommand{\cX}{\mathcal{X}}
\newcommand{\cU}{\mathcal{U}}
\newcommand{\cH}{\mathcal{H}}
\newcommand{\Sh}{\mathfrak{Sh}}
\newcommand{\bF}{\mathbf{F}}
\newcommand{\op}{\textrm{op}}
\newcommand{\fr}{\textrm{fr}}
\newcommand{\SYZ}{\textrm{SYZ}}
\newcommand{\cpx}{\textrm{cpx}}
\newcommand{\Kah}{\textrm{K\"ah}}
\newcommand{\stab}{\textrm{stab}}
\newcommand{\Aut}{\mathrm{Aut}}
\newcommand{\aut}{\mathfrak{aut}}
\newcommand{\Hom}{\mathrm{Hom}}
\newcommand{\rank}{\mathrm{rank}}
\newcommand{\Dom}{\mathrm{Dom}}
\newcommand{\End}{\mathrm{End}}
\newcommand{\Fuk}{\mathrm{Fuk}}
\newcommand{\Ker}{\mathrm{Ker}\,}
\renewcommand{\Im}{\mathrm{Im}\,}
\newcommand{\GL}{\mathrm{GL}}
\newcommand{\SL}{\mathrm{SL}}
\newcommand{\SU}{\mathrm{SU}}
\newcommand{\Sp}{\mathrm{Sp}}
\newcommand{\cE}{\mathcal{E}}
\newcommand{\Jac}{\mathrm{Jac}}
\newcommand{\IMM}{\mathrm{IMM}}

\newcommand{\Span}{\mathrm{Span}}
\newcommand{\conn}{\nabla}
\newcommand{\Lieder}{\mathscr{L}}
\newcommand{\Laplacian}{\Delta}
\newcommand{\der}{\mathrm{d}}
\newcommand{\dd}[1]{\frac{\partial}{\partial #1}}
\newcommand{\dbar}{\bar{\partial}}
\newcommand{\contract}[2]{\iota_{#1} #2}
\newcommand{\dotprod}[2]{\left\langle #1 \, , \, #2 \right\rangle}
\newcommand{\pairing}[2]{\left( #1\, , \, #2 \right)}
\newcommand{\conj}[1]{\overline{#1}}
\newcommand{\WT}[1]{\widetilde{#1}}
\newcommand{\OL}[1]{\overline{#1}}
\newcommand{\oi}{\overline{\iota}}
\newcommand{\tr}{\mathrm{tr} \,}
\newcommand{\unu}{\underline{\nu}}
\newcommand{\Mor}{\mathrm{Mor}}
\newcommand{\nat}{\mathbb{N}}
\newcommand{\integer}{\mathbb{Z}}
\newcommand{\rat}{\mathbb{Q}}
\newcommand{\real}{\mathbb{R}}
\newcommand{\quat}{\mathbb{H}}
\newcommand{\oct}{\mathbb{O}}
\newcommand{\one}{\mathbf{1}}
\newcommand{\consti}{\mathbf{i}\,}
\newcommand{\conste}{\mathbf{e}}
\newcommand{\ev}{\mathrm{ev}}
\newcommand{\sphere}[1]{\mathbf{S}^{#1}}
\newcommand{\proj}{\mathbb{P}}
\newcommand{\bP}{\mathbb{P}}
\newcommand{\torus}{\mathbf{T}}
\newcommand{\moduli}{\mathcal{M}}
\newcommand{\bM}{\mathcal{M}}
\newcommand{\bMC}{\mathcal{MC}}
\newcommand{\cA}{\mathcal{A}}
\newcommand{\cB}{\mathcal{B}}
\newcommand{\cC}{\mathcal{C}}
\newcommand{\cD}{\mathcal{D}}
\newcommand{\cO}{\mathcal{O}}
\newcommand{\cI}{\mathcal{I}}
\newcommand{\cJ}{\mathcal{J}}
\newcommand{\cL}{\mathcal{L}}
\newcommand{\cM}{\mathcal{M}}
\newcommand{\pt}{\mathrm{pt}}
\newcommand{\AI}{A_\infty}
\newcommand{\MF}{\mathrm{MF}}
\newcommand{\Tw}{\mathrm{Tw}}
\newcommand{\LM}{\mathcal{LM}}
\newcommand{\QH}{\mathrm{QH}}
\newcommand{\CM}{\mathcal{LF}}
\newcommand{\mirror}{\mathrm{mirror}}
\newcommand{\barL}{\bar{L}}
\newcommand{\new}{\mathrm{new}}
\newcommand{\DT}{\mathrm{DT}}
\newcommand{\Coh}{\mathrm{Coh}}
\newcommand{\val}{\mathrm{val}}
\newcommand{\cd}{\check{\partial}}
\newcommand{\Id}{\mathrm{Id}}
\newcommand{\dt}[1]{\doubletilde{$#1$}}
\newcommand{\st}[1]{\tilde{#1}}
\newcommand{\lmod}{{\operatorname{-mod}}}
\newcommand{\tL}{\tilde{L}}
\newcommand{\diag}{\mathrm{diag}}
\newcommand{\cG}{\mathscr{G}}


\DeclareMathOperator*{\argmin}{arg\,min}
\DeclareMathOperator*{\argmax}{arg\,max}




%%
%% work harder in optimizing text layout. Typically shrinks text by 1/6 of page, enable
%% it at the very end of the writing process, when you are just above the page limit
%%
% \usepackage{microtype}

%%
%% fine-tune paragraph spacing
%%
% \renewcommand{\paragraph}[1]{\vspace{.5em}\noindent\textbf{#1.}}

%%
%% globally adjusts space between figure and caption
%%
% \setlength{\abovecaptionskip}{.5em}


%%
%% Allows "the use of \paper to refer to the project name"
%% with automatic management of space at the end of the word
%%
% \usepackage{xspace}
% \newcommand{\paper}{ProjectName\xspace}


%%
%% Tigthen underline
%%
% \usepackage{soul}
% \setuldepth{foobar}


% Todonotes is useful during development; simply uncomment the next line
%    and comment out the line below the next line to turn off comments
%\usepackage[disable,textsize=tiny]{todonotes}
\usepackage[textsize=tiny]{todonotes}

% The \icmltitle you define below is probably too long as a header.
% Therefore, a short form for the running title is supplied here:
%\icmltitlerunning{Submission and Formatting Instructions for ICML 2026}

\begin{document}

\twocolumn[
  \icmltitle{Laws of Learning Dynamics and the Core of Learners}

  % It is OKAY to include author information, even for blind submissions: the
  % style file will automatically remove it for you unless you've provided
  % the [accepted] option to the icml2026 package.

  % List of affiliations: The first argument should be a (short) identifier you
  % will use later to specify author affiliations Academic affiliations
  % should list Department, University, City, Region, Country Industry
  % affiliations should list Company, City, Region, Country

  % You can specify symbols, otherwise they are numbered in order. Ideally, you
  % should not use this facility. Affiliations will be numbered in order of
  % appearance and this is the preferred way.
  \icmlsetsymbol{equal}{*}

  \begin{icmlauthorlist}
    \icmlauthor{Inkee Jung}{equal,BU}
    \icmlauthor{Siu Cheong Lau}{equal,BU}
    %\icmlauthor{}{sch}
    %\icmlauthor{}{sch}
  \end{icmlauthorlist}

  \icmlaffiliation{BU}{Department of mathematics, Boston University, Boston, USA}

  \icmlcorrespondingauthor{Inkee Jung}{inkeej@bu.edu}
  \icmlcorrespondingauthor{Siu Cheong Lau}{scllouis@bu.edu}

  % You may provide any keywords that you find helpful for describing your
  % paper; these are used to populate the "keywords" metadata in the PDF but
  % will not be shown in the document
  

  \icmlkeywords{Learning Dynamics, Adversarial Attack, Deep Neural Network, Ensemble Machine Learning, Logifold}

  \vskip 0.3in
]

% Note on paperlimit 
%%%%%%%%%%%%%%%%

% Page limit: The main body of the paper has to be fitted to 8 pages,
% excluding references and appendices; the space for the latter two is not
% limited in pages, but the total file size may not exceed 10MB.

%%%%%%%%%%%%%%%%%

% this must go after the closing bracket ] following \twocolumn[ ...

% This command actually creates the footnote in the first column listing the
% affiliations and the copyright notice. The command takes one argument, which
% is text to display at the start of the footnote. The \icmlEqualContribution
% command is standard text for equal contribution. Remove it (just {}) if you
% do not need this facility.

% Use ONE of the following lines. DO NOT remove the command.
% If you have no special notice, KEEP empty braces:
%\printAffiliationsAndNotice{}  % no special notice (required even if empty)
% Or, if applicable, use the standard equal contribution text:
\printAffiliationsAndNotice{\icmlEqualContribution}


\begin{abstract}
	We formulate the fundamental laws governing learning dynamics, namely the conservation law and the decrease of total entropy. Within this framework, we introduce an entropy-based lifelong ensemble learning method. 
	We evaluate its effectiveness by constructing an immunization mechanism to defend against transfer-based adversarial attacks on the CIFAR-10 dataset. 
	Compared with a naive ensemble formed by simply averaging models specialized on clean and adversarial samples, the resulting logifold achieves higher accuracy in most test cases, with particularly large gains under strong perturbations.
\end{abstract}

\section{Introduction}
In this paper, we develop a rigorous mathematical formulation of learning process of ensembles of models, and prove two fundamental laws governing learning dynamics: the conservation law and the decrease of entropy, which bear a strong formal resemblance to the First and Third Laws of Thermodynamics. Learning dynamics is a very rich and foundational subject, see for instance \cite{Markov_decision_processes, NNLearning,Nonlinear_systems, Dynamic_programming,  loss_surfaces, saddle,SGD_Bayesian,  mean_field, bias, convergence_generalization}. We will take a concise, mathematical and practical approach.

Entropy and cross entropy play crucial roles in machine learning and data science, providing a quantitative framework for measuring uncertainty and loss in predictive modeling, see for instance \cite{Hendrycks2017, Guo2017,decomposition_uncertainty}. The concept of entropy has a long history, originating in 19th-century thermodynamics. The formulation most directly relevant to data science is Shannon entropy, introduced in his mathematical theory of communication \cite{Shannon}. 

A model $(\Phi,f)$ for the truth function $\cT$ is defined as a composition of two functions: the Euclidean representation $\Phi$ that maps a part of the full domain to an Euclidean space (called a feature space), and the decision function $f$ that maps the feature space to a subset of the truth labels. 
We study a modular AI architecture \cite{society_of_mind,neural_module} in which an ensemble is composed of heterogeneous models acting on different input modalities and output subspaces. This puts modular ensembles, Mixture-of-Experts \cite{adaptive_MOE} and multimodal learning \cite{multimodal,multimodal_Boltzmann,multimodal_survey} into a unified framework, which have the advantages of reduced catastrophic forgetting \cite{Google_nested}, improved accuracy \cite{conditional_computation,Sparsely_Gated_MOE,scaling_giant_conditional,switch_transformers,transferable_visual_language}, and finer control over training through targeted fine-tuning and model-level error attribution \cite{ST_MoE}.

The entropy of an ensemble, as formulated in this work, consists of two principal components: the total entropy over the union of the domains covered by the constituent models, and a measure of the complement of this union within the full domain of the underlying ground-truth function.
This formulation reveals two fundamental ingredients of the learning process. The first is reinforcement and refinement \cite{stat_grad_RL,RL_book,  policy_grad_RL, human_deep_RL, trust_policy_optimization, ProximalPO, soft_actor_critic} within the union of domains already covered by existing models. The second is the expansion of the knowledge domain \cite{learning_domains,domain_inv_feature,domain_adv_train,lost_domain} through the training of new models on newly observed or generated data.


We implement this theory of learning process to locate untrustworthy predictions and to develop a defense mechanism against common types of adversarial attacks \cite{AdvGoodfellow,AdvCW,APGD,auto_attack}. The idea is to construct a logifold \cite{Jung-Lau2,Jung-Lau} that consists of two layers of models (also called generations). The first layer contains two types of models: those trained on original samples and those trained on weakly perturbed samples. The second layer consists of models trained on incoming (possibly strongly) perturbed samples detected by the first layer. We use the total entropy of the first layer to identify perturbed samples which are then passed to the second layer. Beyond the static rejection or abstention mechanisms proposed in prior work \cite{CARL, Pang2022, CPR}, our system treats rejected samples as active signals to evolve into a newly immunized generation.


In this design, the first layer functions as an immunized system capable of detecting adversarial perturbations via entropy. Once a non-negligible volume of inputs is detected outside the core, the system transfers models from the first layer to a second layer, where additional training is performed using the detected adversarial samples. This hierarchical structure can be extended to multiple layers to strengthen defenses; however, in our experiments we restrict attention to a two-layer architecture.
We evaluate the proposed approach on the CIFAR-10 \cite{CIFAR-10} dataset under gradient-based attacks such as APGD and AutoAttack \cite{auto_attack, APGD}. With immunization, our logifold achieves significantly higher accuracy than models trained solely on clean data, models trained solely on adversarially perturbed data, and a naive ensemble obtained by simply averaging all models.

\section{Laws of Learning for a model}
Let 
$$\cT:X \to Y$$ 
be a continuous measurable function that represents the truth where $X$ is a topological measure space with $\mu(X) < +\infty$ and $1<|Y|<\infty$. In other words, we assume that the actual world we want to model has no fuzziness.

\begin{definition} \label{def:model}
	A model $(\Phi,f)$ for the truth function $\cT$ is a continuous measurable map $\Phi: S \to \R^n$ for a measurable subset $S \subset X$ and a function 
	$$f=(f_1,\ldots,f_t): \R^n \to \Delta(T)$$ 
	which is continuous on the image of $\Phi$, and $\Delta(T)$ is the simplex of probability distributions on a finite subset $T \subset Y$ of $t$ elements with $T \supset \cT(S)$. $\Phi$ is called to be an Euclidean representation and $f$ is called to be a decision function.
\end{definition}

\begin{definition}[Entropy of a model] \label{def:entropy}
	Let
	\begin{equation}
		H:\Delta(T) \to \R_{\geq 0}, \,\,H(y) = -\sum_{a=1}^t y_a \log_t y_a \label{eq:H(y)}
	\end{equation} 
	be the Shannon entropy function.
	The total entropy of a model $(\Phi,f)$ for $\cT$ is 
	$$ H(\Phi,f) := \int_S H(f_a(\Phi(x))) dx + \mu(X-\Dom(\Phi)). $$
\end{definition}
The second term comes from integration over $X - \Dom(\Phi)$ of the maximal entropy value which is $\sum_{a=1}^t (-\frac{1}{|Y|} \log_{|Y|} \frac{1}{|Y|}) = 1$. 

It is well-understood that the entropy quantifies the uncertainty of the model $f$. In other words,
\begin{center}
	{\bf Entropy is fuzziness.}
\end{center} 
Note that $H(\Phi,f)$ does not depend on the truth function $\cT$: we do not need to know about $\cT$ in order to compute the entropy $H(\Phi,f)$.

Next, we consider the distance between $f$ and the truth function $\cT$ which quantifies the error of the model $f$. 

\begin{definition}[Cross entropy between the truth and a model] \label{def:cross}
	Let
	$$ h: \Delta \times \Delta^\circ \to \R_{\geq 0},\,\, h(y,\gamma) = -\sum_{a=1}^d y_a \log \gamma_a $$
	be the cross entropy function, where $\Delta = \{y\in \R_{\geq 0}^d: \sum_{a=1}^d y_a = 1\}$ and $\Delta^\circ$ denotes its interior.
	The cross entropy between the truth function $\cT$ and a model $(\Phi,f)$ for $\cT$ is defined to be $H(\cT,(\Phi,f)) = \int_S h(\cT(x), f(\Phi(x))) \,dx + \mu(X-S)$, which equals
	$$ H(\cT,(\Phi,f)) := \int_S -\log_t \left(f_{\cT(x)}(\Phi(x))\right)  dx + \mu(X-S)$$
	where $f_{y}(s) := 0$ if $y \not\in T$.
\end{definition}
The second term comes from the integration over $X-S$ of the maximal cross entropy value $-\log_{|Y|} \frac{1}{|Y|} = 1$.
The first term is $+\infty$ if the measure $\mu(\{x \in S: \cT(x) \not\in T\}) > 0$. We have avoided this by putting the requirement $\cT(S) \subset T$ in Definition \ref{def:model}.

The difference between cross entropy and entropy is quantified by the following proposition. The left hand side is the total work done needed to go from $(\Phi,f)$ to $\cT$ against the gradient vector field of $H$ that measures fuzziness. Interpreting the cross entropy $H(-,(\Phi,f))$ as potential, then the right hand side is the potential difference between $\cT$ and $(\Phi,f)$. Thus, the identity is an analog of {\bf the conservation law (The First Law of Thermodynamics)}: 
\begin{center}
	{\bf Effort paid in learning the truth $\cT$ turns to potential.}
\end{center}

\begin{proposition}[Conservation Law of Learning] \label{prop:PD}
\begin{align*}
&\int_{\Dom(\Phi)} (\mathrm{grad}\, H)|_{f(\Phi(x))} \cdot (\cT(x) - f(\Phi(x))) \,dx\\
=& H(\cT,(\Phi,f)) - H(\Phi,f).
\end{align*}
\end{proposition}
\begin{proof}
	See Appendix~\ref{app:proof:prop:PD}.
\end{proof}

\begin{definition}
	A learning process of a model for $\cT$ is a sequence of models $(\Phi^{(i)},f^{(i)})_{i=1}^N$ for $N \in \Z_{>1} \cup \{+\infty\}$ such that $\mathrm{Dom}(\Phi^{(i)}) \subset \mathrm{Dom}(\Phi^{(j)})$ and $H(\cT,(\Phi^{(i)},f^{(i)})) > H(\cT,(\Phi^{(j)},f^{(j)}))$ for $i<j$.
\end{definition}

The \textbf{Law of Learning Dynamics} as:
\begin{center}
	\textbf{Entropy decreases in a learning process.}
\end{center}
This has interesting similarity with the Third Law of Thermodynamics for crystallization.

\begin{theorem} \label{thm:law_single}
	Suppose that $(\Phi^{(i)},f^{(i)})_{i=1}^\infty$ is a learning process such that $\lim_{i\to\infty} H(\cT,(\Phi^{(i)},f^{(i)})) = 0$. Then 
	$$\lim_{i\to\infty} H(\Phi^{(i)},f^{(i)}) = 0.$$
	In particular, there exists a learning subprocess $(\Phi^{(i_j)},f^{(i_j)})_{j=1}^\infty$ such that the entropy decreases to zero.
\end{theorem}

\begin{proof}
	First, note that both $-\log_{t_i}\left(f^{(i)}_{\cT(x)}(\Phi^{(i)}(x))\right)$ and $-f^{(i)}(\Phi^{(i)}(x)) \log_{t_i} f^{(i)}(\Phi^{(i)}(x))$ are non-negative functions. By $\lim_{i\to\infty} H(\cT,(\Phi^{(i)},f^{(i)})) = 0$, for any $\epsilon, n>0$, we can take $i$ sufficiently large such that $\mu(X-\mathrm{Dom}(\Phi^{(i)})) < \epsilon$ and 
	$$ \int_{\mathrm{Dom}(\Phi^{(i)})} -\log_{t_i} \left(f^{(i)}_{\cT(x)}(\Phi^{(i)}(x))\right) < \frac{\epsilon}{n}$$
	so that $$\mu\left(\left\{-\log_{t_i}\left(f^{(i)}_{\cT(x)}(\Phi^{(i)}(x))\right) > \frac{1}{n}\right\}\right) < \epsilon.$$
	We need to prove that the following quantity can be made arbitrary small:
	\begin{align*}
		& \int_{\mathrm{Dom}(\Phi^{(i)})} \sum_{a=1}^{t_i} \left(-f^{(i)}_a(\Phi^{(i)}(x)) \log_{t_i} f^{(i)}_a(\Phi^{(i)}(x)) \right) dx \\
		%=& \sum_{a,b=1}^{t_i} \int_{\mathrm{Dom}(\Phi^{(i)})\cap \cT^{-1}\{b\}}  \left(-f^{(i)}_a(\Phi^{(i)}(x)) \log_{t_i} f^{(i)}_a(\Phi^{(i)}(x)) \right) dx\\
		%=& \sum_{a=b} \int_{\mathrm{Dom}(\Phi^{(i)})\cap \cT^{-1}\{a\}}  \left(-f^{(i)}_a(\Phi^{(i)}(x)) \log_{t_i} f^{(i)}_a(\Phi^{(i)}(x)) \right) dx\\
		%+& \sum_{a\not=b} \int_{\mathrm{Dom}(\Phi^{(i)})\cap \cT^{-1}\{b\}}  \left(-f^{(i)}_a(\Phi^{(i)}(x)) \log_{t_i} f^{(i)}_a(\Phi^{(i)}(x)) \right) dx\\
		=& \int_{\mathrm{Dom}(\Phi^{(i)})} \left(-f^{(i)}_{\cT(x)}(\Phi^{(i)}(x)) \log_{t_i} f^{(i)}_{\cT(x)}(\Phi^{(i)}(x)) \right) dx\\
		+& \int_{\mathrm{Dom}(\Phi^{(i)})} \sum_{a\not=\cT(x)} \left(-f^{(i)}_a(\Phi^{(i)}(x)) \log_{t_i} f^{(i)}_a(\Phi^{(i)}(x)) \right) dx.\\
	\end{align*}
	For the first term, we subdivide $\mathrm{Dom}(\Phi^{(i)}) = A \cup B$ into two subsets according to whether $-\log_{t_i} f^{(i)}_{\cT(x)}(\Phi^{(i)}(x)) \leq \frac{1}{n}$ or not. On $A$, the integrand $-f^{(i)}_{\cT(x)}(\Phi^{(i)}(x)) \log_{t_i} f^{(i)}_{\cT(x)}(\Phi^{(i)}(x)) \leq \frac{t_i^{-1/n}}{n} < \frac{1}{n}$, and hence the integral over $A$ is less than $\frac{\mu(X)}{n} < \epsilon$ by taking $n$ sufficiently large.  On $B$, since $\mu(B) < \epsilon$ and the integrand $\leq \frac{1}{t_i}$, the integral is less than $\epsilon$.
	
	For the second term, we note that the set $C = \{x \in \mathrm{Dom}(\Phi^{(i)}):f^{(i)}_a(\Phi^{(i)}(x)) > 1 - t_i^{-1/n} \textrm{ for some } a\not=\cT(x)\}$ is a subset of $B$ which has measure $< \epsilon$. Since the integrand $< 1$, the integral over $C$ is less than $\epsilon$. In the complement of $C$, the integrand $\leq -(t_i-1)(1 - t_i^{-1/n}) \log_{t_i} (1 - t_i^{-1/n}) \leq -(|Y|-1)(1 - |Y|^{-1/n}) \log_{2} (1 - |Y|^{-1/n})$ which can be controlled to be $<\epsilon / \mu(X)$ by taking $n$ large enough, so that the integral is less than $\epsilon$. This finishes the proof.
\end{proof}

\section{Law of Learning for an ensemble}
\begin{definition}
An ensemble is a collection of models $\cU = \{(\Phi_l,f_l): l=1,\ldots,K\}$. Its knowledge domain is the union $\Dom(\cU) = \bigcup_{l=1}^k \Dom(\Phi_l)$.
\end{definition}

\begin{definition}[Entropy for ensemble] \label{def:H(U)}
For two models $(\Phi_1,f_1)$ and $(\Phi_2,f_2)$, consider $F_i := f_i \circ \Phi_i: \Dom(\Phi_1) \cap \Dom(\Phi_2) \to \Delta(T_1 \cup T_2)$. The cross entropy over their common intersection is defined to be
$ H((\Phi_1,f_1),(\Phi_2,f_2)) := \int_{\Dom(\Phi_1) \cap \Dom(\Phi_2)} h(F_{1}(x),F_{2}(x)) \,dx. $

For an ensemble of models $\cU = \{(\Phi_l,f_l): l=1,\ldots,K\}$, 
the pointwise entropy of $\cU$ at $x \in X$ is defined as
\begin{equation} \label{eq:H_x(U)}
	H_x(\cU) = \frac{1}{K_x^2} \sum_{\substack{l: \Dom(\Phi_l)\ni x \\\\r: \Dom(\Phi_r)\ni x}} h(F_{l}(x),F_{r}(x))
\end{equation}
if $x \in \Dom(\cU)$ and $K_x>0$ is the number of models with $\Dom(\Phi_l) \ni x$; otherwise $H_x(\cU)$ is
$$\frac{1}{K^2} \sum_{l,r=1}^K \sum_{a=1}^{|Y|}-\frac{1}{|Y|}\log_{|Y|}\frac{1}{|Y|} = 1.$$
The total entropy is defined to be
\begin{equation} \label{eq:H(U)}
	H(\cU) = \int_X H_x(\cU) = \int_{\Dom(\cU)} H_x(\cU) + \mu(X-\Dom(\cU)).
\end{equation}


For two ensembles of models $\cU = \{(\Phi_l,f_l): l=1,\ldots,L\}, \mathcal{V} = \{(\Psi_r,g_r): r=1,\ldots,R\}$, the pointwise cross entropy of $(\cU,\mathcal{V})$ at $x \in X$ is defined as

\begin{equation} \label{eq:H_x(U,V)}
	H_x(\cU,\mathcal{V}) = \frac{1}{L_{x} R_{x}} \sum_{\substack{l: \Dom(\Phi_l)\ni x \\\\r: \Dom(\Psi_r)\ni x}} h((f_l\circ \Phi_{l})(x),(g_r\circ \Psi_r)(x))
\end{equation}
if $L_x >0$ and $R_x >0$ where $L_x = |\{l: \Dom(\Phi_l)\ni x\}|$ and $R_x = |\{r: \Dom(\Psi_r)\ni x\}|$; 
$$ \frac{-1}{R_x |Y|} \sum_{r: \Dom(\Psi_r)\ni x} \log_{t_r} \prod_{a=1}^{t_r} (g_{r,a}\circ \Psi_{r})(x) $$
if $L_x = 0$ and $R_x > 0$, and
$$ \frac{-1}{L_x} \sum_{l: \Dom(\Phi_l)\ni x} \sum_a F_{l,a} \log_{|Y|} \frac{1}{|Y|} = 1$$
otherwise. The total cross entropy $H(\cU,\mathcal{V})$ is the integration of $H_x(\cU,\mathcal{V})$ over $X$.

In particular, the pointwise cross entropy of the truth function $\cT$ and $\cU$ is the average cross entropy:
\begin{equation} \label{eq:H_x(T,U)}
	H_x(\cT,\cU) = \frac{1}{L_{x}} \sum_{l: \Dom(\Phi_l)\ni x} h(\cT(x),(f_l\circ \Phi_{l})(x))
\end{equation}
if $x \in \Dom(\cU)$ and $1$ otherwise.
\end{definition}

\begin{remark} \label{rmk:entropy_average}
	We can also consider the entropy $H(\bar{F})$ of the average function $\bar{F}$ of $\{(\Phi_i,f_i)\}_{i=1}^K$, which is defined by $\bar{F}(x) = \frac{1}{L}\sum_{l: \Dom(\Phi_l) \ni x} f_i(\Phi_i(x))$ if $L := |\{l: \Dom(\Phi_l) \ni x\}|>0$ and $\bar{F}_a(x) = 1/|Y|$ for all $a$ otherwise. This is different from the entropy of the ensemble $H(\cU)$ defined above and loses information about disagreements of models in the ensemble. For instance, suppose $\cU$ consists of two model functions $F_1,F_2$ with $F_1(x_1)=(0.6,0.4), F_2(x_1) = (0.6,0.4)$ at $x_1$, and $F_1(x_2)=(0.95,0.05), F_2(x_2) = (0.25,0.75)$ at $x_2$. At the two points, the average is the same, $\bar{F}(x_1) = \bar{F}(x_2) = (0.6,0.4)$, and hence $H_{x_1}(\bar{F}) = H_{x_2}(\bar{F})$. On the other hand, since $F_1$ and $F_2$ have big disagreement at $x_2$, $H_{x_2}(\cU) > H_{x_1}(\cU)$.
\end{remark}

\begin{remark}
For a broad information-theoretic analysis of predictive-uncertainty measures, including cross-entropy based formulations, we refer the reader to \cite{schweighofer2025}.
\end{remark}

$H(\cU)$ is the degree of fuzziness of the ensemble of models $\cU$. In particular,
\begin{theorem}[Vanishing of entropy]\label{thm:van_ent}
For an ensemble of models $\cU = \{(\Phi_l,f_l): l=1,\ldots,K\}$, 
$$H(\cU)=0$$ 
if and only if the models cover the whole space in measure sense, that is $X - \bigcup_{l=1}^K \Dom(\Phi_l)$ has measure zero, and there exists a function $F: X \to Y$ such that for every $l = 1,\ldots,K$, $F|_{\Dom(\Phi_l)} = f_l \circ \Phi_l$ almost everywhere. 
\end{theorem}

\begin{proof}
	See Appendix~\ref{app:proof:thm:van_ent}.
\end{proof}

When each model in the ensemble is given by a fuzzy linear logical function (which gives a mathematical interpretation of a neural network), this is called to be a \textbf{fuzzy linear logifold}, see Definition 9 of \cite{Jung-Lau}. A strict linear logifold is given by Definition 8 of \cite{Jung-Lau}. The terminology of a logifold is motivated from manifold theory in which the space is covered by charts of observers. 
For a manifold, different charts are required to agree with each other and are glued together over overlapping regions; in contrast, fuzziness and disagreement among charts are allowed for a fuzzy logifold and are measured by entropy.
By the above theorem,
\begin{corollary}
	A fuzzy linear logifold is a strict linear logifold if and only if its entropy vanishes.
\end{corollary}

\begin{definition} \label{def:core}
Given $0 < p < \min\{\frac{2}{|Y|}, \log_{|Y|} \frac{3}{2}\}$, the core of an ensemble of models $\cU = \{(\Phi_l,f_l): l=1,\ldots,K\}$ in threshold $p$ is defined to be the subset
$$ C_p(\cU) = \left\{x \in X: H_x(\cU) < p \right\} \subset X $$
where $H_x(\cU)$ is the pointwise entropy in Definition \ref{def:H(U)}.
\end{definition}

\begin{proposition}[Non-fuzzy limit of the core]\label{prop:NonFuzLimCore}
	In the core $C_p(\cU)$ of an ensemble $\cU = \{(\Phi_l,f_l): l=1,\ldots,K\}$ where $0 < p < \min\{\frac{2}{|Y|}, \log_{|Y|} \frac{3}{2}\}$, for every $x\in C_p(\cU)$ and $l$ such that $\Dom(\Phi_l) \ni x$, there exists a unique $a=a_{l,x} \in T_l$ that $f_{l,a}\circ\Phi_l(x)\in [0,1]$ achieves maximum; moreover, $a_{l,x}$ are the same for all such $l$. Thus, this defines a function 
	$$F: C_p(\cU) \to Y$$ 
	which is called to be the non-fuzzy limit of the ensemble restricted to the core.
\end{proposition}

\begin{proof}
	See Appendix~\ref{app:proof:prop:NonFuzLimCore}.
\end{proof}

Below is an analog of Proposition \ref{prop:PD} for ensemble learning. The proof is similar and is omitted. The left hand side is the work done to go from the average of the ensemble to the truth $\cT$ against the average gradient vectors of $H$ that measures fuzziness. This turns to the potential energy defined by the cross entropy $H(-,\cU)$. 

\begin{proposition}[Conservation law of ensemble learning]
\begin{equation}
\int_{\Dom(\cU)} V(x) \cdot (\cT(x) - f(x)) \,dx = H(\cT,\cU) - H(\cU)
\end{equation}
where 
\begin{equation}
	f(x) = \frac{1}{L_x}\sum_{l: \Dom(\Phi_l)\ni x} f_l(\Phi_l(x))
\end{equation}
and
\begin{equation}
	V(x) = \frac{1}{L_x}\sum_{l: \Dom(\Phi_l)\ni x} (\mathrm{grad}\, H)|_{f_l(\Phi(x))}
\end{equation}
denote the average of models and the average of gradient vectors respectively.
\end{proposition}
 
\begin{definition}[Learning process of an ensemble] \label{def:learning_ensemble}
	A learning process of an ensemble of models for $\cT$ is a sequence of ensembles $(\cU^{(i)})_{i=1}^N$ and injective functions 
	$$\mathcal{L}:\cU^{(i)} \to \cU^{(i+1)}$$ 
	for $i=1,\ldots,N-1$ such that $\Dom(\cU^{(i)}) \subset \Dom(\cU^{(j)})$ and $H(\cT,\cU^{(i)}) > H(\cT,\cU^{(j)})$ for $i < j$, with the additional assumption that if $N=\infty$, the domain and the target of each model stabilizes: for each $i$ and each element $\Phi \in \cU^{(i)}$, there exists $k$ such that $\Dom(\mathcal{L}^j(\Phi)) = \Dom(\mathcal{L}^{j+1}(\Phi))$ and $T(\mathcal{L}^j(\Phi)) = T(\mathcal{L}^{j+1}(\Phi))$ for all $j\geq k$. The models in $\cU^{(i+1)}$ that are not in the image of $\mathcal{L}_i$ are said to be created in the $i$-th step; the models $\Phi \in \cU^{(i+1)}$ with $\Dom(\Phi) = \emptyset$ are said to be annihilated in the $i$-th step.
\end{definition}

We have the following \textbf{Law of Learning Dynamics for ensemble}:

\begin{theorem}\label{thm:LLDforEns}
Let $(\cU^{(i)})_{i=1}^\infty$ be a learning process of an ensemble of models for $\cT$ such that $\lim_{i\to\infty} H(\cT,\cU^{(i)}) = 0$. Then 
$$\lim_{i\to\infty} H(\cU^{(i)}) = 0.$$
In particular, there exists a learning subprocess $(\cU^{(i_j)})_{j=1}^\infty$ such that the entropy $H(\cU^{(i_j)})$ decreases to zero.
\end{theorem}

\begin{proof}
	See Appendix~\ref{app:proof:thm:LLDforEns}.
\end{proof}


%This paper focuses on detection of mistakes made by models. To know what a model does not know.  We set two thresholds of cross entropy of a logifold, low and high.  The predictions with low entropy are highly reliable.  This is the core of the logifold.  The predictions with high entropy are very likely to be wrong.  (That is, P(wrongly predicted|high entropy) is high.)  

%We develop an ensemble-based approach to predict samples that will be wrongly predicted. With a system of multiple classifiers, we quantify the degree of disagreement using cross entropy of their outputs.  We apply this to detect and defend adversarial attacks.  The successful rate of predicting wrongly predicted samples with high entropy is XX\%.  Moreover, for attacks by small perturbation, the accuracy of the system is XX\%, which is XX\% gain compared to the original single model.  

%For adversarial attack with bigger perturbation, which was less studied in the literature, even the ensemble system fails to defend and the accuracy drops to XX\%.  To handle this, we introduce a novel approach based on logifold which is inspired by manifold theory.  We introduce new models that are specialized to classify adversarial samples, and RESTRICT THE DOMAINS of both the original and new models.  The system dynamically allocates the input to the appropriate chart-specific classifiers according to the cross entropy of their outputs. In this situation, our method achieves up to a 30\% improvement in accuracy compared to the original ensemble system.


\section{Entropy-based lifelong ensemble learning}

Lifelong learning addresses situations where a learner faces a stream of learning tasks, providing an opportunity for knowledge transfer across tasks \cite{Thrun1998, overcoming_forgetting, memory_based_lifelong, lifelong_review}.

In the context of the learning process in Definition \ref{def:learning_ensemble}, the sequence of tasks is to learn the truth function on an expanding subsets of the domain over time. We consider a sequence of multiple learners where the earlier learners retain the memory of previous tasks, and new learners are created and transferred to new tasks. Earlier learners decide whether to pass a given task to the next generation by the entropy which measures uncertainty and disagreement among them. This multi-generational approach complements standard Deep Ensembles \cite{DeepEnsembles, Ovadia2019}, which can be fragile on outlier predictions. Concretely, this is determined by entropy thresholds and the core in Definition \ref{def:core}; the earlier learners are annihilated if their core becomes an empty set.

\begin{definition}[Lifelong learning]
	Let $E$ be a set whose elements are called environments, $A$ be a set whose elements are called actions, and
	$\cX_e = (\cX,\mu_e)$ where $\cX$ is a topological space called the data space with a probability measure $\mu_e$ that depends on environments $e\in E$. A learner is a tuple $(K,U,\pi)$ where $K$ is a topological space called the knowledge state space, $U: K \times \cX \to K$ is called the learning rule, and $\pi:K\times E \to A$ is called the performance policy.
	Let $P: A\times E\to\R$ that evaluates actions in different environments and $I:A \to \R$ that measures memory formation.
	
	A lifelong learning process is an infinite sequence $(e_i,x_i,k_i) \in E \times \cX \times K$, where $x_i$ are drawn according to the measure $\mu_{e_i}$ and $k_i$ is determined by $k_{i} = U(k_{i-1},x_{i-1})$ and the initial $k_0$ such that the evaluation and memory formation are non-decreasing up to a prefixed $\epsilon>0$:
	\begin{align*}
		P(\pi(k_{i+1},e_{i+1}),e_{i+1}) & > P(\pi(k_{i},e_{i+1}),e_{i+1}) - \epsilon; \\
		I(\pi(k_{i+1},e_{i+1})) &> I(\pi(k_{i},e_{i+1})) - \epsilon. 
	\end{align*}
\end{definition}

In a classification task, $\cX$ is the product $X\times Y$ of the input space and the labeling set. In practice, $\cX_e$ is replaced by its own product $\cX_e^N$ for some $N$, meaning that we take a batch of training and validation samples according to the measure $\mu_e$ rather than a single sample in each learning step. For a neural network model, $K$ is the parameter space of the model, $E$ is the set of all possible testing batches and $A$ is the set of all possible predictions (in the form of probability distributions on the labeling set) for the input testing batches.  
$P$ measures the accuracy and $I$ measures the entropy.

However, if the environment changes rapidly and drastically over time, catastrophic forgetting may happen \cite{catastrophic,van_de_Ven_2025} which crashes lifelong learning, namely the inequality for $I$ is violated. For this purpose, we consider an ensemble of learners instead to better retain memory. See also Nested Learning recently introduced in \cite{Google_nested}. 

\begin{definition}[Lifelong ensemble learning]
	\label{def:lifelong_ensemble}
	Let $E$, $A$, $\cX_e$, $P$ be the environment set, the action set, the data space and the evaluation of actions as before. Moreover, $I$ is given as a function $A^n \to \R$ that measures overall memory formation for an arbitrary number $n$ of learners.
	Given an infinite sequence of environments $e_i \in E \times \cX$ which produce samples $x_i \in \cX_{e_i}$,
	a lifelong ensemble learning process is an infinite sequence 
	$$(\cU^{(i)},\Gamma^{(i)}, (k^{(i)}_l:l \in \cU^{(i)}))_{i=1}^\infty$$ 
	where each $\cU^{(i)}$ is an ensemble of finitely many learners, $\Gamma^{(i)}: A^{|\cU^{(i)}|} \to A$ is the aggregation of individual actions (voting), with injective functions $\cL^{(i)}: \cU^{(i)}\to \cU^{(i+1)}$ where $l\in \cU^{(i)}$ and $\cL^{(i)}(l)$ are the same learner (meaning $(K,U,\pi)$ remains the same), and individual knowledge states $k^{(i)}_l\in K^{(i)}_l$ are updated according to $U^{(i)}_l$ as before. The evaluation $P$ and memory formation $I$ for the aggregated action $\Gamma^{(i)}((\pi^{(i)}_l(k_l^{(i)},e_i): l \in \cU^{(i)}))$ are required to satisfy the same inequalities as before.
\end{definition}

The idea is to use entropy thresholds set by accuracy evaluated on validation samples to decide the aggregation policy $\Gamma^{(i)}$. This suggests the following procedure for a classification task.

\begin{enumerate}
	\item We start with a subset of the domain $S^{(1)}\subset X$. Using the data $\{(x,\cT(x)): x\in S^{(1)}\}$ (which is partitioned into training and validation subsets), we train and validate an ensemble of models $\cU^{(1)}$ whose domains are subsets of $S^{(1)}$ that cover $S^{(1)}$ and whose targets are subsets of $\cT(S^{(1)})$. The entropy in Definition \ref{def:H(U)} is measured with respect to $S^{(1)}$. (At this stage, $\mu(X-S^{(1)})$ is dropped since it is unknown and cannot be measured).
	\item For a preset high accuracy threshold, we use the validation data to find a low enough entropy threshold $p$ such that on the core $C_p(\cU^{(1)})$ of Definition \ref{def:core}, the ensemble $\cU^{(1)}$ reaches the required accuracy on the validation data. We restrict the domain of $\cU^{(1)}$ to be $C_p(\cU^{(1)})$, and train the next generation of models $\cU^{(2)'}$ whose domain is restricted to be $S^{(1)} - C_p(\cU^{(1)})$. Then we take $\cU^{(2)} = \cU^{(1)} \cup \cU^{(2)'}$. Note that the entropy threshold $p$ could be zero, in which case the required accuracy threshold could not be reached and the first generation $\cU^{(1)}$ is annihilated.
	\item The above process can be repeated several times until the overall accuracy reaches a satisfactory level. If necessary, we may enlarge the models in the next generation to reach higher accuracy.
	\item Next, we enlarge the domain subset $S^{(1)}$ to $S^{(2)} \subset X$ by finding more samples from the real world or generating more samples by different means such as adding random noises, acting by symmetries, or adversarial attacks. The entropy in Definition \ref{def:H(U)} is updated over the bigger subset $S^{(2)}$ accordingly. The entropy thresholds in each layer of the ensembles are recalculated to match the required accuracy threshold. More generations of models are created as before and some earlier generations are annihilated, until the overall accuracy reaches a satisfactory level.
	\item The above procedure is repeated until the domain $X$ is saturated by training and validation samples.
\end{enumerate}

\section{Immunized systems against adversarial attacks}
\label{sec:experiments}

We empirically support the proposed learning dynamics on the CIFAR-10 dataset.
We construct a system that consists of generations of neural network models, where the domain of each generation is restricted according to the entropies of generations.  
More precisely, the domain of the first generation is restricted to its core (Definition \ref{def:core}); 
the domain of the second generation is restricted to the intersection of its core with the out-core (the complement of the core) of the first generation;
and the domain of the last generation is restricted to the complement of all restricted domains of the previous generation.  
Such an ensemble of neural network models with restricted domains is called a logifold in \cite{Jung-Lau2,Jung-Lau}. This perspective is closely related to selective classification with an abstain/reject option, where uncertain or unsafe inputs are rejected rather than force-classified \cite{CARL, SDIM, CPR}.

Distance-based detectors can be sensitive to preprocessing and other sources of mismatch \cite{Mahalanobis}. We will use \emph{ensemble entropy to detect the change of environments in lifelong learning}. Namely, when an unusual amount of inputs lie outside the core of the logifold, the system determines that the environment has changed and it initiates an adaptation procedure, which resets the entropy thresholds of the generations to reach a satisfactory accuracy for the training samples in the new environment, possibly annihilates some of the early generations if their thresholds are almost zero, and transfers models in the early generations to form a new generation to learn the samples that are out of core. 

In this paper, we use adversarial perturbations primarily as a controlled mechanism to expand the observed domain and to stress test the stability of the logifold.
First, we start with the clean environment $e_0$, where we have a distribution supported over clean samples and train an ensemble of models $\cU^{(0)}$ on these clean samples.  Then we raise a white-box adversarial attack on $\cU^{(0)}$. In the experiment to be described in the next section (see Table \ref{tab:entropy_failure}), we find that the core of $\cU^{(0)}$ can detect the attacks without the knowledge of accuracy rate. Thus, an environmental change from $e_0$ to $e_1$ is detected, and some of the models are transferred to $e_1$ and added to $\cU^{(0)}$ to form a new ensemble $\cU^{(1)}$.

We will see that the core of $\cU^{(1)}$ detects white-box strong adversarial attacks against itself even better than that of $\cU^{(0)}$. For this, we take an even stronger white-box adversarial attacks (in terms of the size of perturbations and the number of steps) against $\cU^{(1)}$. 
Compared to the same degree of attacks to $\cU^{(0)}$, the core coverage of $\cU^{(1)}$ exhibits a more significant drop and the core accuracy of $\cU^{(1)}$ is higher. We call this to be an \textbf{immunization effect}, and call $\cU^{(1)}$ to be an immunized generation whose core is more sensitive to adversarial environments. 

Then the immunization mechanism works as follows. When a non-negligible volume of incoming samples lie outside the core of  $\cU^{(1)}$, the system determines that it is under a strong adversarial environment. Then it adds the second generation $\cU^{(2)^\prime}$ consisting of models that are trained on union of the strong adversarial and original samples that lie outside the core of $\cU^{(1)}$. The final immunized system $\cU^{(2)}$ (which is also denoted as IMM) consists of both generations $\cU^{(1)}$ and $\cU^{(2)^\prime}$. In the testing/application stage, testing samples that lie in the core of $\cU^{(1)}$ are processed by $\cU^{(1)}$, and those that lie outside the core of $\cU^{(1)}$ are processed by $\cU^{(2)^\prime}$. Comparing to a simple average of all models, this mechanism gains a higher accuracy on the whole union of samples, see Table \ref{tab:acc_eval}.

Note that we are not claiming for theoretical robustness concerning worst-case guarantees over all small perturbations of data points, which requires the knowledge of almost all data points and has a trade-off with accuracy \cite{robustness_accuracy, trade_off_robustness}. In \cite{mixing_classifiers}, accurate models are mixed with robust models to  alleviate the accuracy-robustness trade-off. In our design, in addition to taking average of an ensemble, we aim at detecting effective attacks on $\cU^{(1)}$ by using the entropy of $\cU^{(1)}$, and maximizing the accuracy on the union of the original and adversarial samples by the use of entropy of the multi-generation design.


\subsection{Experimental setup}
\label{subsec:exp_setup}

We use the standard CIFAR-10 split (50{,}000 train / 10{,}000 test) and an ensemble of $K=8$ independently trained classifiers: four VGG-type~\cite{VGG} and four ResNet-type~\cite{ResNet} networks.
The baseline ensemble predictor $\bar{f}$ is the probability average.

We evaluate on clean data and on adversarial counterparts generated from the same splits, as well as their union
\[
\mathcal{D}_{\mathrm{union}} \coloneqq 
\mathcal{D}_{\mathrm{clean}} \cup \mathcal{D}_{\mathrm{weak}} \cup \mathcal{D}_{\mathrm{strong}}.
\]
All adversarial examples are $\ell_2$-bounded and clipped to $[0,1]$. 
We run untargeted $\ell_2$ APGD-CE~\cite{APGD} against $\bar f$ with configurations
\[
(\epsilon,\epsilon_{\text{step}},\text{steps},\text{restarts})
\in \{(0.5,0.2,2,4),\ (0.7,0.2,10,4)\},
\]
denoted as \texttt{APGD-weak} ($\cD_{\mathrm{weak}}$) and \texttt{APGD-strong} ($\cD_{\mathrm{strong}}$), respectively.
We also generate transfer-based adversarial examples using \texttt{AutoAttack}~\cite{auto_attack} with surrogate models from RobustBench~\cite{RobustBench}; we use the CIFAR-10 \texttt{Standard} $\ell_2$ surrogate and run AutoAttack at $\epsilon=0.5$.

We replace the integrals in Equation~\eqref{eq:H(U)} for total entropy $H$ by empirical averages over the evaluation set, and drop the complement term as in Definition~\ref{def:lifelong_ensemble}.

\subsection{Specialization, composing generations, and evaluation}
\label{subsec:specialization_and_evaluation}

In the logifold viewpoint, each model acts as a chart whose effective domain and decision boundary may differ across charts.
Rather than training a new network per chart, we reuse the penultimate feature backbone and adapt only the classification head.

\emph{Specialization.} Given a trained classifier and a dataset, we replace the final dense$+$softmax head and fine-tune with Adam~\cite{Adam} and categorical cross-entropy for 21 epochs using a stepwise learning-rate schedule.

\paragraph{Compute environment.}
Experiments were run on the BU SCC cluster using an NVIDIA L40S GPU.


\subsection{Results}
\label{subsec:results}

\paragraph{Ensemble entropy versus single-model entropy.}
We evaluate core coverage and core accuracy under white-box $\ell_2$ APGD-CE (weak/strong) at a fixed entropy threshold $\tau=0.1$.
For a single model, entropy thresholding does not produce a reliable core under adversarial perturbations: many misclassified adversarial samples still receive high confidence, yielding small $H_x$ while core accuracy remains low.


Let $\cU^{(0)}$ be the baseline ensemble (four VGG-type and four ResNet-type models trained on clean data).
Under a white-box attack on $\bar f$, the core of $\cU^{(0)}$ shrinks substantially on adversarial inputs.
This shows that ensemble entropy of Definition \ref{def:H(U)} behaves better than the entropy of the average of \(\cU^{(0)}\) treated as a single model (see Remark \ref{rmk:entropy_average}) for this purpose.

We then add four ResNet specialists trained on \texttt{APGD-weak} samples generated against $\cU^{(0)}$.
Let \texttt{SPs} denote these models and define $\cU^{(1)}=\cU^{(0)}\cup\texttt{SPs}$ ($K_1=12$), which we call the immunized generation.
Table~\ref{tab:entropy_failure} reports that $\cU^{(1)}$ yields a smaller core and higher core accuracy than $\cU^{(0)}$ and \texttt{SPs} under strong attacks.

\begin{table}[t]
	\centering
		\vspace{0.1in}
	\caption{Core coverage and core accuracy under white-box $\ell_2$ APGD-CE attacks (weak/strong configurations) at entropy threshold $\tau=0.95$, which is selected for $U^{(0)}$ to have greater than $90\%$ core coverage on clean sample.
		We report single models (ResNet/VGG), the baseline ensemble $U^{(0)}$, specialists (\texttt{SPs}), and the immunized generation $U^{(1)}$.
		Under strong attacks, single models fail to detect perturbed samples in terms of core coverage. $U^{(1)}$ yields the smallest core and the highest core accuracy among the three ensembles, which best distinguishes the effective adversarial samples among these cases.}
	\label{tab:entropy_failure}
		\vspace{0.1in}
	\setlength{\tabcolsep}{6pt}

	{\scriptsize
		\begin{tabular}{lcccc}
			\toprule
			& &Clean & \makecell{ \texttt{APGD-CE} \\Weak}& \makecell{ \texttt{APGD-CE} \\ strong}\\
			\midrule
			ResNets &\makecell{Cov. (\%) \\ Acc. (\%)}&
			\makecell[c]{ 100  \\  91.35 ($\pm$ 0.67)} &
			\makecell[c]{ 100\\  20.75($\pm$ 4.04)} &
			\makecell[c]{ 100  \\  17.43  ($\pm$ 4.07)} \\
			
			VGGs &\makecell{Cov. (\%) \\ Acc. (\%)}&
			\makecell[c]{ 100\\  91.38($\pm$ 0.82)} &
			\makecell[c]{ 100\\  52.54  ($\pm$ 3.56)} &
			\makecell[c]{100 \\ 46.97($\pm$ 4.94)} \\
			
			$U^{(0)}$&\makecell{Cov. (\%) \\ Acc. (\%)}&
			\makecell[c]{90.50\\ 97.26} &
			\makecell[c]{72.96\\ 88.50 } &
			\makecell[c]{51.07 \\ 22.62 } \\
			
			\texttt{SPs}&\makecell{Cov. (\%) \\ Acc. (\%)}&
			\makecell[c]{94.02 \\ 93.70} &
			\makecell[c]{83.77 \\ 88.70 } &
			\makecell[c]{49.48 \\ 32.32 } \\
			
			$U^{(1)}$&\makecell{Cov. (\%) \\ Acc. (\%)}&
			\makecell[c]{89.38 \\ 97.39 } &
			\makecell[c]{75.08\\ 94.21 } &
			\makecell[c]{43.06\\ 34.58 } \\
			\bottomrule
	\end{tabular}}
	
\end{table}

\paragraph{Evaluation across environments and memory formation.}
We run lifelong ensemble learning over $\cD_{\mathrm{clean}}, \cD_{\mathrm{weak}}, \cD_{\mathrm{strong}}$, and $\cD_{\mathrm{union}}$, where ``weak''/``strong'' refer to adversarial samples generated against $\cU^{(0)}$. 
For each domain, we select an entropy threshold $\tau$ on a validation set to ensure high core accuracy for $\cU^{(1)}$ (see Fig.~\ref{fig:th_sweep} in Appendix~\ref{app:graphs}).


On clean data, $\cC_\tau(\cU^{(1)})$ covers most samples, so we do not construct a second generation.
On the other domains, the out-of-core region becomes non-negligible, and we train out-of-core specialists $\cU^{(2)'}$ and form $\IMM=\cU^{(1)}\cup\cU^{(2)'}$, a two-generation logifold (IMM: \emph{Immunization Mechanism}).

The $\IMM$ aggregator is
\begin{equation}
	\Gamma^{(2)}(\IMM)\coloneqq 
	\mathbf{1}_{\cC_\tau(\cU^{(1)})}\Gamma^{(1)}(\cU^{(1)})
	+\mathbf{1}_{\cC_\tau^c(\cU^{(1)})}\Gamma^{(2)'}(\cU^{(2)'}),
\end{equation}
where $\Gamma^{(1)}$ and $\Gamma^{(2)'}$ are probability averages on their respective ensembles.
We compare overall accuracy $P$ and total entropy $H$ across $\Gamma^{(0)}$ (baseline $\cU^{(0)}$), $\Gamma^{(1)}$ (first generation $\cU^{(1)}$), and $\Gamma^{(2)}$ ($\IMM$) in Table~\ref{tab:eval_summary}.
We also report transfer-based \texttt{AutoAttack} results.

\textbf{Remark.}
In Definition~\ref{def:lifelong_ensemble}, the memory formation functional $I$ is required to be non-decreasing.
In our experiments, we use the total entropy $H$ as a proxy for $I$, but $H$ is expected to decrease as learning progresses.
Thus, when reporting $H$ we implicitly adopt a monotone reparameterization of memory formation. The distinction is a matter of directionality.

\begin{table}[t]
	\centering
	\vspace{0.1in}
	\caption{Overall accuracy and total entropy at entropy thresholds $\tau$ selected on validation data to ensure high first-generation core accuracy. $\cU^{(0)}$: 8 scratch-trained models. $\cU^{(1)}$: 12-model immunized generation. $\IMM=\cU^{(1)}\cup\cU^{(2)'}$ with four out-of-core specialists. ``All-ens'' averages all members in $\IMM$.}
	
	\label{tab:eval_summary}
		\vspace{0.1in}

	
	\begin{subtable}{1.0\linewidth}
		\centering
			\vspace{0.1in}
			\caption{Overall accuracy on each domain at the selected $\tau$. Except for the clean domain, where core coverage is already high, the logifold $\IMM$ performs best.}
			
		\label{tab:acc_eval}
			\vspace{0.1in}
		{\footnotesize
		\begin{tabular}{lcccc}
			\toprule
			\  Dataset & $\cU^{(0)}$ & $\cU^{(1)}$  & All-ens& $\IMM$  \\
			\midrule
			$\cD_{\mathrm{clean}}$, \scriptsize{$\tau = 1.4$}  & \textbf{94.33\%} & 94.23\%& N/A& N/A \\
			$\cD_{\mathrm{weak}}$, \scriptsize{$\tau = 0.96$}& 66.69\%& 77.32\%&81.40\%  & \textbf{85.44\%}\\
			$\cD_{\mathrm{strong}}$, \scriptsize{$\tau = 0.02$} & 13.16\%& 21.48\%&45.50\% & \textbf{90.93\%} \\
			$\cD_{\mathrm{union}}$, \scriptsize{$\tau = 0.66$}      & 58.06\%& 64.34\%& 74.62\%&\textbf{87.48\%} \\  
			\bottomrule
		\end{tabular}}
	\end{subtable}
	
	\begin{subtable}{0.98\linewidth}
		\centering
			\vspace{0.1in}
			\caption{Total entropy $H$ (Def.~\ref{def:H(U)}) on each domain at the selected $\tau$. We observe a decrease in total entropy, which is an important signal of successful lifelong ensemble learning.}
		
		\label{tab:ent_eval}
			\vspace{0.1in}
		\begin{tabular}{lcc}
			\toprule
			\  Dataset & $\cU^{(1)}$  & $\IMM$  \\
			\midrule
			$\cD_{\mathrm{weak}}, \tau = 0.96$& \makecell{$H = 0.6408 $}& \makecell{$H=0.2876 $}
%			$ 55\%$
			\\
			$\cD_{\mathrm{strong}}, \tau = 0.02$ & \makecell{$H =  1.9782$}& \makecell{$H=0.2054$}
%			$ 90\%$
\\
			$\cD_{\mathrm{union}}, \tau = 0.66$      &  \makecell{$H =  0.9638$}& \makecell{$H=0.2317$}
%			$75.96\%$
\\  
			\bottomrule
		\end{tabular}
\end{subtable}
\end{table}

Table~\ref{tab:acc_eval} shows that IMM improves overall accuracy on $\cD_{\mathrm{union}}$ (from $64.34\%$ with $\cU^{(1)}$ to $87.48\%$), and also yields the best performance on $\cD_{\mathrm{weak}}$ and $\cD_{\mathrm{strong}}$ among the compared aggregation rules.
Given changing environments, we observe that the inequality in lifelong ensemble learning $P(\Gamma^{(1)}) < P(\Gamma^{(2)})$ holds.
 Table~\ref{tab:ent_eval} reports a consistent reduction in total entropy from $\cU^{(1)}$ to IMM on non-clean environments.

Finally, we consider transfer-based \texttt{AutoAttack} at $\epsilon=0.5$ and evaluate $U^{(0)},U^{(1)}$ and $\IMM$ evolved on $\cD_{\mathrm{union}}$ with $\tau=0.66$.
The accuracy of $\cU^{(0)}, \cU^{(1)}, \IMM$ are $78.47\%$, $84.71\%$ and $83.58\%$ respectively. The result is comparable with $\cD_{\mathrm{weak}}$ shown in Table \ref{tab:acc_eval}, which suggests that the \texttt{AutoAttack} under consideration is slightly weaker than $\cD_{\mathrm{weak}}$ and hence $\cU^{(1)}$ attains the highest overall accuracy. On the other hand, we observe that even in this case, $\IMM$ increases core coverage from $0.4581$ ($\cU^{(1)}$) to $0.7096$ while improving high core accuracy from $95.87\%$ ($\cU^{(1)}$) to $97.04\%$.


\section{Conclusion}
We established two fundamental laws concerning work done and entropy for learning dynamics of models and ensembles. We formulated the notion of cores for ensembles, which utilizes the ensemble entropy as a label-free quantity of fuzziness and disagreement.

We constructed an immunized architecture IMM for detecting and adapting to polluted environments by adversarial perturbations. IMM consists of an immunized generation consisting of both models trained on clean data and models transferred to weakly perturbed data to detect effective adversarial samples, and a specialized generation that consists of models transferred to out-of-core samples. On CIFAR-10 under APGD-CE and surrogate-based \texttt{AutoAttack}, the resulting system IMM improves accuracy on mixed domains and yields particularly large gains under strong perturbations compared to naive averaging baselines.
%Our specialization is intentionally targeted to ``polluted'' samples under a controlled mixture ratio. When inputs are known to be unpolluted, we set the entropy threshold sufficiently high so that predictions are produced only by the first layer.

More broadly, our approach suggests a principled way to grow modular ensembles over expanding domains by (i) identifying the core (trustworthy regions) via low ensemble entropy, (ii) transferring models to the complement of the core, and (iii) iterating this process to enlarge the core.


\section*{Impact Statement}

This paper presents work whose goal is to advance the field of Machine
Learning. There are many potential societal consequences of our work, none
which we feel must be specifically highlighted here.

% In the unusual situation where you want a paper to appear in the
% references without citing it in the main text, use \nocite
%\nocite{langley00}x


\bibliographystyle{icml2026}
\bibliography{bib}
%%%%%%%%%%%%%%%%%%%%%%%%%%%%%%%%%%%%%%%%%%%%%%%%%%%%%%%%%%%%%%%%%%%%%%%%%%%%%%%
%%%%%%%%%%%%%%%%%%%%%%%%%%%%%%%%%%%%%%%%%%%%%%%%%%%%%%%%%%%%%%%%%%%%%%%%%%%%%%%
% APPENDIX
%%%%%%%%%%%%%%%%%%%%%%%%%%%%%%%%%%%%%%%%%%%%%%%%%%%%%%%%%%%%%%%%%%%%%%%%%%%%%%%
%%%%%%%%%%%%%%%%%%%%%%%%%%%%%%%%%%%%%%%%%%%%%%%%%%%%%%%%%%%%%%%%%%%%%%%%%%%%%%%
\newpage
\appendix
\onecolumn

\section{Proofs}

\subsection{Proof of Proposition~\ref{prop:PD}}
\label{app:proof:prop:PD}
\begin{proof}
	$$-(\mathrm{grad}\, H)|_{y} \cdot y' = -h(y',y) + \sum_{a} y'_a $$
	by direct calculation. for $y'$ being tangent to $\Delta$, the second term vanishes. Thus the integrand on the left hand side becomes
	$$ h(\cT(x)-F(x),F(x)) = h(\cT(x),F(x)) - H(F(x))$$
	where $F = f \circ \Phi$.
	Integration over $\Dom(\Phi)$ gives the equality.
\end{proof}

\subsection{Proof of Theorem~\ref{thm:van_ent}}
\label{app:proof:thm:van_ent}
\begin{proof}
	
	Note that each term of $H(\cU)$ in Definition \ref{def:H(U)} is non-negative. Thus, $H(\cU)=0$ if and only if $\mu(X - \bigcup_{l=1}^K \Dom(\Phi_l))=0$ and $H((\Phi_l,f_l),(\Phi_r,f_r))=0$ for every $l,r$. $H((\Phi_l,f_l),(\Phi_r,f_r))$ is an integral of a non-negative function. By monotone convergence theorem in measure theory, the integral is zero if and only if the non-negative functions $-f_{l,a} (\Phi_l(x)) \log_t f_{r,a} (\Phi_r(x)) = 0$ almost everywhere. For $l=r$, this happens if and only if $f_{l,a}(\Phi_l(x)) \in \{0,1\}$ for $x \in S_l'$ where $\mu(\Dom(\Phi_l) - S_l')=0$; since the image of $f_l$ lies in the simplex $\Delta(T_l)$, $f_{l,a}(\Phi_l(x)) = 1$ for exactly one value of $a \in T_l$ and $0$ otherwise; hence it gives a function $F_l$ on $S_l'$. For $l\not=r$, this happens if and only if $f_{l,a}(\Phi_l(x)) = f_{r,a}(\Phi_r(x))$ for $x \in S_{lr}'$ where $\mu(\Dom(\Phi_l) \cap \Dom(\Phi_r) - S')=0$. In this situation, for each $x \in X$, we define $F(x) = F_1(x)$ if $x \in S_1'$, $F(x) = F_2(x)$ if $x \in S_2'-S_1'$, $F(x) = F_3(x)$ if $x \in S_3'-S_2'-S_1'$ and so on; $F(x)$ is assigned to be a fixed element in $Y$ if $x \in X-\bigcup_i S_i'$ which is a measure zero set. By construction, $F|_{\Dom(\Phi_l)} = f_l \circ \Phi_l$ almost everywhere for every $l = 1,\ldots,K$. 
\end{proof}

\subsection{Proof of Proposition~\ref{prop:NonFuzLimCore}}
\label{app:proof:prop:NonFuzLimCore}
\begin{proof}
	
	First, we have $C_p(\cU) \subset \bigcup_l \Dom(\Phi_l)$ since $H_x(\cU) < \log_{|Y|} 2 < K^2$ for all $x\in C_p(\cU)$. Moreover, since every term in \eqref{eq:H_x(U)} is non-negative, $H_x(\cU) < \log_{|Y|} 2$ implies that every term $H_x((\Phi_l,f_l),(\Phi_r,f_r)) = \sum_{a=1}^t -F_{l,a}(x) \,\log_t F_{r,a}(x) < \log_{|Y|} 2$ where $F_i = f_i \circ \Phi_i$ and $t = |T_l \cup T_r|$. Let's consider the terms for $l=r$. The subset $\{y \in \Delta(T_l): \sum_{a=1}^t -y_a \,\log_t y_a < \log_{|Y|} 2\}$ has $t$ connected components; for all $y$ lying in the $a$-th component, $y_a$ is the unique maximum among the entries of $y$. This proves the first statement for the uniqueness of $a \in T_l$. Now consider $y_l = F_{l}(x)$ and $y_r = F_{r}(x)$ for $x\in C_p(\cU)$ and $\Dom(\Phi_l) \cap \Dom(\Phi_r) \ni x$.  For every $a \in T$, it needs to satisfy the inequalities 
	\begin{align*}
		-(y_{l,a} \,\log_t y_{r,a} + y_{r,a} \,\log_t y_{l,a}) < 2/|Y| <& 2/t; \\
		-(y_{l,a} \,\log_t y_{r,a} + y_{r,a} \,\log_t y_{l,a}) < \log_{|Y|} \frac{3}{2} <& \log_t \frac{3}{2}
	\end{align*}
	where the second inequality reduces to
	$$-(y_{l,a} \,\log y_{r,a} + y_{r,a} \,\log y_{l,a})< \log \frac{3}{2}.$$ 
	For each $a$, the subset of all $(y_{l,a},y_{r,a})\in [0,1]\times [0,1]$ that satisfies the inequalities consists of two components that are subsets of $\{\textrm{both } y_{l,a} \textrm{ and } y_{r,a} < \frac{1}{|Y|}\}$ (using the first inequality) and $\{\textrm{both } y_{l,a} \textrm{ and } y_{r,a} > \frac{1}{2}\}$ (using the second inequality) respectively. For those indices $a$ that are not argmax of $y_l$, we have $y_{l,a} < 1/2$ which implies $y_{r,a} < \frac{1}{|Y|} < \frac{1}{|T_r|}$. Let  $a_0$ be the argmax of $y_l$. Then $y_{r,a_0} = 1 - \sum_{a\not=a_0} y_{r,a} > \frac{1}{|T_r|}>\frac{1}{|Y|}$, and hence $(y_{l,a_0},y_{r,a_0})$ lies in the second component which implies $a_0$ is also the argmax of $y_r$. Thus $y_l$ and $y_r$ have the same argmax $a_0$ Thus models give the same response (after taking argmax) in the intersections, and hence we have a well-defined function $F:C_p(\cU)\to Y$.
\end{proof}


\subsection{Proof of Theorem~\ref{thm:LLDforEns}}
\label{app:proof:thm:LLDforEns}
\begin{proof}
	
	By \eqref{eq:H_x(T,U)} and Theorem \ref{thm:law_single}, the entropy of each model in an ensemble that survives in the limit has entropy limiting to zero. Thus, it suffices to consider the cross entropy of any two different models in an ensemble. Since each domain and target stabilize, we may assume that their domains and targets are independent of the step $i$. Let's denote their domains by $D_l \subset X$, targets by $T_l \subset Y$ and the model functions by $F^{(i)}_l= (F^{(i)}_{l,a}:a\in T_l)$ respectively for $l=1,2$. Take any $\epsilon > 0$and $n>0$. For $i$ large enough, as in the proof of Theorem \ref{thm:law_single}, we have a subset $S_l \subset D_l$ with $\mu(D_l - S_l) < \epsilon$ such that for all $x \in S$,
	$$ -\log_{t_l} F_{l,\cT(x)}^{(i)}(x) < \frac{1}{n}, $$
	that is, $F_{l,\cT(x)}^{(i)}(x) > 1/(t_l)^{1/n}$ (meaning that it can be made arbitrarily close to $1$). The pointwise cross entropy between the two is
	$$ -F_{1,\cT(x)}^{(i)}(x) \log_{t_2} F_{2,\cT(x)}^{(i)}(x) - \sum_{a\not=\cT} F_{1,a}^{(i)}(x) \log_{t_2} F_{2,a}^{(i)}(x). $$
	Let $x \in S_1 \cap S_2$. For the first term, $-\log_{t_2} F_{2,\cT(x)}^{(i)}(x) < 1/n$ and $F_{1,\cT(x)}^{(i)}(x) < 1$, implying that the first term $<1/n$. For each summand of the second term, $F_{l,a}^{(i)}(x) < 1-1/(t_l)^{1/n} < 1-1/|Y|^{1/n} < \epsilon / |Y|$ by taking $n$ large. Thus, the pointwise entropy for $x \in S_1 \cap S_2$ can be made arbitrarily small. Moreover, $\mu(D_1 \cap D_2 - S_1 \cap S_2) < \epsilon$. Thus, the cross entropy between any two models in $\cU^{(i)}$ can be made arbitrarily small for $i$ sufficiently large. Thus, $H(\cU^{(i)})$ defined by \eqref{eq:H(U)} can be made arbitrarily small.
\end{proof}
\clearpage
\section{Graphs of Core Coverage and Accuracy under Entropy thresholding.}
\label{app:graphs}
\begin{figure}[!htbp]
	\centering
	
	\begin{subfigure}[t]{0.49\linewidth}
		\centering
		
		\includegraphics[width=\linewidth]{./figures/union_original_weak_baseline_strong_baseline_val_counts_and_acc_dualaxis.png}
		\vspace{0.1in}
		\caption{$\mathcal{D}_{\mathrm{union, val}}$. At $\tau  = 0.7$ the core accuracy is greater than $95\%$.}
		\label{fig:th_sweep_union}
		\vspace{0.1in}
	\end{subfigure}\hfill
	\begin{subfigure}[t]{0.49\linewidth}
		\centering
		\includegraphics[width=\linewidth]{./figures/original_val_counts_and_acc_dualaxis.png}
		\vspace{0.1in}
		\caption{$\mathcal{D}_{\mathrm{clean, val}}$. At $\tau = 1.4$, the core coverage is 0.9704.}
		\label{fig:th_sweep_clean}
		\vspace{0.1in}
	\end{subfigure}
	
	\begin{subfigure}[t]{0.32\linewidth}
		\centering
		\includegraphics[width=\linewidth]{./figures/APGD_strong_baseline_val_counts_and_acc_dualaxis.png}
		\vspace{0.1in}
		\caption{$\mathcal{D}_{\mathrm{strong, val}}, \tau = 0.02$}
		\vspace{0.1in}
		\label{fig:th_sweep_APGD_strong}
	\end{subfigure}\hfill
	\begin{subfigure}[t]{0.32\linewidth}
		\centering
		\includegraphics[width=\linewidth]{./figures/APGD_weak_baseline_val_counts_and_acc_dualaxis.png}
		\vspace{0.1in}
		\caption{$\mathcal{D}_{\mathrm{weak}}, \tau = 0.96$}
		\label{fig:th_sweep_APGD_weak}
		\vspace{0.1in}
	\end{subfigure}\hfill
	\begin{subfigure}[t]{0.32\linewidth}
		\centering
		\includegraphics[width=\linewidth]{./figures/AutoAttack_val_counts_and_acc_dualaxis.png}
		\vspace{0.1in}
		\caption{\texttt{AutoAttack} ($\epsilon=0.5$). The vertical line is at $\tau = 0.1$. It shows that the immunized generation detects transfer-based adversarial samples reliably.}
		\label{fig:th_sweep_aa_std}
		\vspace{0.1in}
	\end{subfigure}
	\vspace{0.1in}	
	
	\caption{
		Accuracy and coverage under entropy thresholding across domains.
		Solid curves show accuracies on $\cC_\tau$ (black) and $\cC_\tau^{c}$ (red).
		Dashed curves show coverages of $\cC_\tau$ (orange) and $\cC_\tau^{c}$ (blue).
		The vertical dashed line marks the selected $\tau$ for a validation accuracy target of $95.25\%$. Across domains, decreasing \(\tau\) monotonically decreases core coverage while increasing core accuracy. Moreover, the core accuracy is consistently above the overall accuracy, whereas the out-of-core accuracy is consistently below it.
	}
	\vspace{0.1in}
\label{fig:th_sweep}
\end{figure}

\end{document}

